\title{Direct Preference Optimization: Your Language Model is Secretly a Reward Model}

\begin{abstract}
Although large-scale unsupervised language models (LMs) acquire extensive knowledge of the world and certain reasoning abilities, exerting fine-grained control over their outputs remains challenging, primarily due to their fully unsupervised training processes. Current approaches aiming to improve controllability typically involve collecting human assessments comparing different model outputs and subsequently fine-tuning the original LM to better reflect these human judgments—most commonly utilizing reinforcement learning from human feedback (RLHF). Yet, RLHF poses significant practical difficulties, as it generally involves first training a separate reward model to capture user preferences, and then employing reinforcement learning to optimize the LM for this surrogate reward while maintaining similarity to the pretrained LM. This work proposes a novel way of parameterizing the reward model within RLHF, enabling analytical derivation of the corresponding optimal policy and thus reducing the standard RLHF objective to a straightforward classification loss. This leads to the \textit{Direct Preference Optimization} (DPO) algorithm, which is stable, effective, and resource-efficient: it removes the need for sampling from the LM during the adaptation phase and avoids extensive hyperparameter search. Experimental results indicate that DPO is capable of aligning LMs with human preferences as well as, or even better than, prior methods. Specifically, DPO outperforms PPO-based RLHF at modulating output sentiment, and matches or surpasses state-of-the-art methods in response quality for summarization and single-turn dialogue tasks, all while being significantly easier to implement and train.
\end{abstract}

\section{Introduction}
Large unsupervised language models (LMs) developed on vast text corpora have demonstrated remarkable emergent abilities~\citep{chowdhery2022palm, brown2020language, touvron2023llama,bubeck2023sparks}. Nevertheless, these models are trained on data authored by humans with diverse intentions, backgrounds, and levels of expertise. Not all of these human qualities are desirable for the models to emulate; for instance, although it is useful for an AI programming assistant to \textit{recognize} common coding errors in order to correct them, we ultimately wish the model to preferentially produce code reflective of the (sometimes uncommon) high-caliber programming found within its training distribution. By the same token, while a language model should be \textit{cognizant} of misconceptions prevalent among large fractions of the population, such as a misconception held by 50\% of people, we do not want it to assert this erroneous belief as correct in half of all relevant interactions. Thus, the challenge lies in distinguishing and enforcing the model's \emph{intended behaviors and outputs} from its broader reservoir of \textit{acquired knowledge and skills}, which is vital for ensuring that AI systems are reliable, controllable, and high-performing~\citep{ouyang2022training}. Most prevailing techniques guide LMs to align with human preferences via reinforcement learning (RL); however, we will demonstrate that the RL-based objective currently employed in these systems can actually be achieved directly using a simple binary cross-entropy loss, thereby streamlining the process of preference optimization.

Broadly speaking, current approaches confer desirable behaviors upon language models through carefully collected datasets of human preferences, which serve as exemplars of behavior considered safe and beneficial. This process of preference learning typically follows an initial stage involving large-scale unsupervised pre-training on diverse text data. Although supervised fine-tuning using human-annotated, high-quality response demonstrations is the most direct method for preference learning, the dominant and most effective paradigm is reinforcement learning from human or AI feedback (RLHF/RLAIF; \citep{christiano2017deep,bai2022constitutional}). In RLHF frameworks, a reward model is trained on preference data elicited from humans, and RL is subsequently employed to tune the language model's policy toward producing responses that maximize these rewards, while maintaining proximity to the original model to prevent undesirable divergence. Despite the impressive fluency and technical competence yielded by RLHF-tuned models, the RLHF pipeline is substantially more intricate than conventional supervised learning. It requires training several separate models and executing model-in-the-loop sampling, which considerably increases computational demands.

In this work, we present a method for optimizing language models according to human preferences without resorting to explicit reward modeling or reinforcement learning techniques. Our proposed approach, \textit{Direct Preference Optimization} (DPO), is an algorithm that inherently maximizes the same objective as standard RLHF methods (maximizing rewards under a KL-divergence penalty), yet offers simplicity in implementation and training. Fundamentally, DPO modifies the model by increasing the ratio of the log probabilities assigned to preferred responses relative to less-preferred alternatives, while employing a dynamic, per-example importance weighting scheme. This mechanism counters the model collapse that often arises with a straightforward probability ratio objective. Like prior approaches, DPO draws on a theoretical preference model (such as the Bradley-Terry model; \cite{bradley1952rankanalysis}) to quantify alignment between a given reward function and observed preference judgments. Distinctively, instead of using the preference model to first define a loss for training a reward function, and subsequently optimizing a policy with this reward, DPO leverages a change of variables to formulate the preference loss directly in terms of the policy. When provided with a dataset reflecting human judgments on model outputs, DPO can thus learn the policy by optimizing a straightforward binary cross-entropy loss, resulting in a policy that is optimal with respect to an implicitly defined reward that matches the observed preferences.

The primary contribution of this paper is Direct Preference Optimization (DPO), a straightforward algorithm that dispenses with reinforcement learning in order to train language models from preference data. Our empirical findings demonstrate that DPO achieves performance on par with or exceeding current methods—including PPO-based RLHF—for a variety of preference-based tasks, such as sentiment control, text summarization, and conversation, and remains effective for language models up to 6B parameters in size.

\section{Related Work}

Recent advances in self-supervised language models have enabled zero-shot task completion \citep{radford2019language} and improved few-shot performance with prompt-based approaches \citep{gpt3,megatron,chowdhery2022palm}. Nonetheless, these models achieve superior results on downstream tasks and show greater alignment with user intent when further fine-tuned using datasets comprised of instructions and human-provided completions \citep{mishra-etal-2022-cross,sanh2022multitask,chung2022scaling,thoppilan2022lamda}. This process, known as `instruction-tuning,' has proven effective at broadening LLM generalization beyond the scope of the tuning dataset and enhances their practical usability \citep{chung2022scaling}. While instruction-tuning has yielded promising outcomes, collecting \textit{relative} human assessments of response quality is often more feasible than obtaining expert demonstrations. Accordingly, subsequent research has leveraged datasets annotated with human preferences to further train LLMs, leading to gains in areas such as translation \citep{kreutzer-etal-2018-reliability}, summarization \citep{stiennon2022learning,ziegler2020finetuning}, narrative generation \citep{ziegler2020finetuning}, and following instructions \citep{ouyang2022training,ramamurthy2023is}. These approaches initially fit a neural reward model to the dataset using a preference-learning framework such as the Bradley-Terry model \citep{bradley1952rankanalysis}, followed by reinforcement learning-based fine-tuning—often employing REINFORCE \citep{williams1992reinforce}, proximal policy optimization (PPO; \cite{schulman2017proximal}), or their variants \citep{ramamurthy2023is}—to optimize for the learned rewards. Related efforts utilize LLMs that have been fine-tuned with human feedback for instruction-following to create further synthetic preference datasets focused on aspects like safety or harmlessness \citep{bai2022constitutional}, relying on light human oversight by providing textual rubrics to guide annotation. This body of work effectively bridges the gap between two research trajectories: reinforcement learning-based optimization of language models for diverse targets~\citep{Ranzato2015SequenceLT,paulus2018a,wu2018learning}, and broader strategies for learning from human preference data \citep{christiano2017deep,kupcsik2018learning}. However, despite the advantages of using relative human preferences for tuning, applying reinforcement learning at scale to LLMs is still fraught with practical difficulties; the present study introduces a theoretically sound method for leveraging relative preference signals without relying on RL.

Beyond language-based settings, the problem of learning policies from preferences has been explored within both bandit and reinforcement learning frameworks, leading to a variety of methodological developments. The contextual dueling bandit (CDB; \cite{yue2012karmed, dudik2015contextual}) extends contextual bandit models to the case where action selection is guided by comparative preferences or rankings, rather than explicit reward signals. In settings lacking absolute reward feedback, the CDB literature replaces the standard notion of policy optimality with the concept of a \textit{von Neumann winner}: a policy guaranteed to achieve an expected win rate of at least 50\% when evaluated against any alternative policy \citep{dudik2015contextual}. Unlike CDBs, in which preference data is received in an online fashion, human preference-based policy learning more commonly leverages fixed, offline collections of action pairwise preference annotations \citep{yan2022human}. 

Similarly, \textit{preference-based RL} (PbRL) operates by using binary preferences, obtained from an \textit{unknown} underlying 'scoring' function, as the learning signal rather than direct rewards \citep{BusaFekete2014, ruiz2023dueling}. PbRL encompasses a range of algorithms—some of which support off-policy preference data reuse—but these approaches typically follow a two-step process: first, an explicit estimation of the underlying scoring (reward) function, and second, policy optimization using this model \citep{jain2013learning, BusaFekete2014, christiano2017deep, sadigh2017active, kupcsik2018learning}. Departing from these multi-stage pipelines, we propose a one-stage policy optimization strategy that directly aligns the policy with observed preferences.

\section{Preliminaries}\label{section:prelims}

We begin by summarizing the RLHF procedure as presented in \citeauthor{ziegler2020finetuning} and expanded in later work (\citep{stiennon2022learning, bai2022training, ouyang2022training}). This framework is structured into three distinct components: (1) supervised fine-tuning (SFT), (2) collection of preference data and reward model estimation, and (3) reinforcement learning-based policy refinement.

\textbf{SFT}: The RLHF workflow initiates with supervised fine-tuning of a pre-existing language model using high-quality examples pertaining to the target applications (such as dialogue or summarization tasks), resulting in a supervised model $\pi^\text{SFT}$. 

\textbf{Reward Modelling Phase}: During the subsequent phase, the SFT-trained model is conditioned on input prompts $x$ to generate pairs of candidate responses $(y_1, y_2)\sim \pi^\text{SFT}(y \mid x)$. Human evaluators are then tasked with indicating their preferred completion for each prompt, denoted $y_w\succ y_l \mid x$ (where $y_w$ and $y_l$ refer to the preferred and less preferred completions, respectively, in the sampled pair). Preferences are assumed to reflect judgments generated by an unobserved, latent reward model $r^*(y, x)$. Several probabilistic models have been employed to fit this human preference data; one commonly adopted method is the Bradley-Terry (BT) model \cite{bradley1952rankanalysis}, although more expressive alternatives such as Plackett-Luce ranking models \citep{plackett1975analysis, luce2012individual} can also be employed given access to multi-way ranked outputs. In the BT framework, the likelihood of observed human preferences $p^*$ is expressed as follows:

\begin{equation}\label{eq:bradley-terry}
    p^*(y_1\succ y_2 \mid x)=\frac{\exp\left(r^*(x, y_1)\right)}{\exp\left(r^*(x, y_1)\right) + \exp\left(r^*(x, y_2)\right)}.
\end{equation}

Given access to a fixed dataset of comparative judgments $\mathcal{D}=\bigl\{x^{(i)}, y_w^{(i)}, y_l^{(i)}\bigr\}_{i=1}^N$ sampled according to $p^*$, a reward model can be parameterized as $r_{\phi}(x, y)$. The parameters of this reward function are then learned by maximizing the likelihood of the observed data. Recasting this as a binary classification problem, the objective corresponds to the following negative log-likelihood loss:

\begin{equation}\label{eq:reward_model}
    \mathcal{L}_R(r_{\phi}, \mathcal{D}) = -\mathbb{E}_{(x, y_w, y_l)\sim \mathcal{D}}\bigl[\log \sigma(r_{\phi}(x, y_w)- r_{\phi}(x, y_l))\bigr]
\end{equation}

Here, $\sigma$ denotes the logistic sigmoid function. In large language model settings, $r_{\phi}(x, y)$ is commonly initialized from the SFT model $\pi^\text{SFT}(y \mid x)$ by augmenting it with a linear output layer applied to the last hidden state of the transformer; this scalar serves as the predicted reward \cite{ziegler2020finetuning}. To mitigate reward variance, it is common practice to normalize the predicted rewards such that $\mathbb{E}_{x,y\sim \mathcal{D}}\left[r_\phi(x, y)\right] = 0$ across all $x$.

\textbf{RL Fine-Tuning Phase}: In the reinforcement learning fine-tuning stage, the acquired reward function guides feedback to the language model. In line with previous work~\citep{jaques2017sequence, jaques2020human}, optimization is formalized as:

\begin{equation}\label{eq:RL}
\max_{\pi_{\theta}}  \mathbb{E}_{x\sim \mathcal{D}, y\sim \pi_{\theta}(y \mid x)}\bigl[r_{\phi}(x, y)\bigr] - \beta\mathbb{D}_{\textrm{KL}}\bigl[\pi_{\theta}(y\mid x)\mid \mid \pi_\text{ref}(y\mid x)\bigr],
\end{equation}

with $\beta$ being a scaling factor that controls how much $\pi_\theta$ is permitted to diverge from a reference policy $\pi_\text{ref}$—typically, this is the initial SFT policy $\pi^\text{SFT}$. In practical implementations, $\pi_\theta$ is initialized to match $\pi^\text{SFT}$. Including this constraint helps restrict the policy from shifting outside the data regime where the reward model is valid, ensuring diversity in generations and reducing the risk of mode collapse to singular, high-reward responses. Since the action space in text generation is discrete, this objective is non-differentiable, and is therefore optimized through reinforcement learning methods. The conventional method \citep{ziegler2020finetuning, stiennon2022learning, bai2022training, ouyang2022training} defines the reward as ${r(x, y) = r_{\phi}(x, y) -\beta (\log \pi_{\theta}(y\mid x) - \log \pi_\text{ref}(y\mid x))}$, and applies PPO~\cite{schulman2017proximal} for maximization.

\section{Direct Preference Optimization}\label{sec:DPO}

Recognizing the complexity and inefficiencies associated with reinforcement learning methods for large-scale language model fine-tuning, we aim to develop a straightforward procedure for policy optimization that directly uses preference data. Contrasting with standard RLHF pipelines, which require training a reward model and then optimizing the policy with reinforcement learning, our approach utilizes a particular form of reward model parameterization that permits analytic recovery of the optimal policy, eliminating the need for an RL loop. As elaborated in the next section, the main conceptual contribution is to exploit a closed-form mapping from reward models to their optimal policies, thereby permitting transformation of reward-based loss functions to policy-based objectives. This change-of-variables approach removes the requirement for explicit reward model fitting while

\textbf{DPO Objective Derivation.} The starting point is the standard RL objective from previous research, Eq.~\ref{eq:RL}, considered for a generic reward function $r$. In line with earlier work~\citep{peters2007reinforcement, peng2019advantage, korbak2022reinforcement, go2023aligning}, it follows readily that the solution to the reward maximization problem with a KL constraint (Eq.~\ref{eq:RL}) is given by:

\begin{equation}\label{eq:op_policy}
    \pi_r(y\mid x) = \frac{1}{Z(x)}\pi_\text{ref}(y\mid x)\exp\left(\frac{1}{\beta}r(x, y)\right),
\end{equation}

where $Z(x) =\sum_{y}\pi_\text{ref}(y\mid x)\exp\left(\frac{1}{\beta}r(x, y)\right)$ serves as the normalization factor or partition function. A full step-by-step derivation is provided in Appendix \ref{app:derivation1}. Direct utilization of Eq.~\ref{eq:op_policy} is impeded in practice by the computational cost of calculating the partition function $Z(x)$, even if one substitutes the MLE-estimated reward $r_{\phi}$ for the true reward $r^*$ \citep{korbak2022reinforcement, go2023aligning}. However, one can manipulate Eq.~\ref{eq:op_policy} to express the reward in terms of the induced optimal policy $\pi_r$, the reference policy $\pi_\text{ref}$, and the partition function. Specifically, by applying the logarithm to both sides and rearranging, we have:

\begin{equation}\label{eq:main_eq}
    r(x,y) =\beta \log \frac{\pi_r(y\mid x)}{\pi_\text{ref}(y\mid x)} + \beta \log Z(x).
\end{equation}

This change of variables may also be performed for the true reward $r^*$ and its associated optimal policy $\pi^*$. Importantly, within the Bradley-Terry model, the relevant probability depends exclusively on the reward difference between two outputs: ${p^*(y_1 \succ y_2 \mid x) = \sigma(r^*(x, y_1) - r^*(x, y_2))}$. Substituting the reparameterization from Eq.~\ref{eq:main_eq} for $r^*(x,y)$ into the Bradley-Terry preference expression in Eq.~\ref{eq:bradley-terry}, we see that the partition function $Z(x)$ drops out, yielding a formulation where preference probabilities are described only in terms of $\pi^*$ and $\pi_\text{ref}$. In summary, the Bradley-Terry model for the optimal RLHF policy $\pi^*$ takes the following form:

\begin{equation}\label{eq:objective}
    p^*(y_1\succ y_2 \mid x)=\frac{1}{1 + \exp\left(\beta \log \frac{\pi^*(y_2\mid x)}{\pi_\text{ref}(y_2\mid x)} - \beta \log \frac{\pi^*(y_1\mid x)}{\pi_\text{ref}(y_1\mid x)}\right)}
\end{equation}

The derivation can be found in Appendix~\ref{app:derivation2}. While this result is derived using the Bradley-Terry model, analogous derivations apply under the Plackett-Luce family~\citep{plackett1975analysis, luce2012individual}, as discussed in Appendix~\ref{app:plackett_luce_models}.

With this policy-based representation for the likelihood of human preference data, it becomes possible to optimize a likelihood for a policy parameterized by $\theta$, denoted $\pi_\theta$. This is directly analogous to the reward modeling loss (cf. Eq.~\ref{eq:reward_model}). The objective function for the policy is then:

\begin{equation}\label{eq:optimum_model}
    \mathcal{L}_\text{DPO}(\pi_{\theta}; \pi_\text{ref}) = -\mathbb{E}_{(x, y_w, y_l)\sim \mathcal{D}}\left[\log \sigma \left(\beta \log \frac{\pi_{\theta}(y_w\mid x)}{\pi_\text{ref}(y_w\mid x)} - \beta \log \frac{\pi_{\theta

In this formulation, we model an implicit reward via a different parameterization, where the resulting optimal policy coincides with $\pi_\theta$. Additionally, because this approach aligns with fitting a reparametrized Bradley-Terry model, it inherits certain desirable theoretical properties, including consistency under appropriate assumptions on the distribution of preference data \cite{bong2022generalized}. Section~\ref{sec:theory} expands on these theoretical aspects and relates DPO to existing literature.

\textbf{Mechanistic interpretation of the DPO update.} To gain insight into the functioning of DPO, it is instructive to examine the gradient of the loss $\mathcal{L}_\text{DPO}$. The gradient with respect to $\theta$ is given by:

\begin{multline*}\label{eq:gradient}
    \nabla_\theta \mathcal{L}_\text{DPO}(\pi_\theta;\pi_\text{ref}) = \\ -\beta\mathbb{E}_{(x, y_w, y_l) \sim \mathcal{D}} \left[\underbrace{\sigma(\hat{r}_\theta(x, y_l) - \hat{r}_\theta (x, y_w))}_\text{greater magnitude when reward ranking is incorrect}\left[\underbrace{\nabla_\theta\log \pi(y_w \mid x)}_\text{boosts probability of $y_w$} - \underbrace{\nabla_\theta\log\pi(y_l \mid x)}_\text{reduces probability of $y_l$}\right]\right],
\end{multline*}

with $\hat{r}_\theta(x, y) = \beta \log \frac{\pi_\theta(y \mid x)}{\pi_\text{ref}(y \mid x)}$, which defines the reward implicitly via the language model $\pi_\theta$ in reference to $\pi_\text{ref}$ (see also Section~\ref{sec:theory}). Conceptually, the gradient term for $\mathcal{L}_\text{DPO}$ encourages increasing the likelihood for favored completions $y_w$ while simultaneously suppressing the likelihood for less favored completions $y_l$. Notably, the contribution of each training example is scaled by the extent to which the implicit reward model $\hat{r}_\theta$ overrates the less favored completion, modulated by $\beta$, thereby reflecting both the model’s ranking accuracy and the severity of the KL regularization. Our empirical results indicate that this reweighting is crucial—simply omitting it, as in a basic version of the approach, can result in model collapse (see Appendix Table~\ref{tab:unlikelihood_generations}).

\textbf{DPO overview.} The typical DPO workflow involves the following steps: 1) For each prompt $x$, generate candidate completions $y_1, y_2 \sim \pi_\text{ref}(\cdot \mid x)$, and assign human preference labels to form an offline preference dataset $\mathcal{D} = \{x^{(i)}, y_w^{(i)}, y_l^{(i)}\}_{i=1}^N$; 2) Adjust the language model $\pi_\theta$ to minimize $\mathcal{L}_\text{DPO}$, given $\pi_\text{ref}$, the dataset $\mathcal{D}$, and a specified $\beta$. 
In real-world applications, it is preferable to utilize already available public preference datasets rather than generating new completions and collecting fresh human feedback. As these datasets are typically constructed using $\pi^\text{SFT}$, we set $\pi_\text{ref} = \pi^\text{SFT}$ when such a model is accessible. Conversely, if $\pi^\text{SFT}$ is unavailable, $\pi_\text{ref}$ is initialized by maximizing the likelihood of preferred completions ${(x, y_w)}$, formulated as ${\pi_\text{ref} = \argmax_{\pi}\mathbb{E}_{x, y_w \sim \mathcal{D}}\left[\log \pi(y_w \mid x)\right]}$. This approach alleviates issues arising from discrepancies between the actual, but inaccessible, reference distribution and the $\pi_\text{ref}$ adopted in DPO. Additional information regarding implementation details and hyperparameter selection is provided in Appendix~\ref{app:implementation}.

\section{Theoretical Analysis of DPO}
In this part, we delve deeper into the theoretical understanding of DPO, offering justification for the method and elucidating how DPO addresses shortcomings found in actor-critic approaches commonly used in RLHF (such as PPO~\cite{schulman2017proximal}).

\label{sec:theory}

\subsection{Your Language Model Is Secretly a Reward Model}
DPO uniquely enables the unification of reward learning and policy optimization into a single maximum likelihood framework, foregoing the need for explicit reward model fitting and separate RL steps. Notably, the objective in Eq.~\ref{eq:main_eq} coincides with the Bradley-Terry model, with a reward formulation given by $r^*(x, y) = \beta \log\frac{\pi^*_\theta(y \mid x)}{\pi_\text{ref}(y \mid x)}$, where we directly optimize the parametric model $\pi_\theta$—paralleling the process of optimizing the reward model in Eq.~\ref{eq:reward_model} after a change of variables. In what follows, we present the theoretical foundations of this reparameterization, establish that it does not restrict the expressivity of possible reward models, and demonstrate that it can fully recover the optimal policy. We start by introducing an equivalence relationship among reward functions.

\begin{definition}
We define two reward functions $r(x, y)$ and $r'(x, y)$ as equivalent if and only if ${r(x, y)-r'(x, y) = f(x)}$ for some function $f$.
\end{definition}

This definition clearly describes an equivalence relation, thereby dividing the space of reward functions into equivalence classes. We now present the following lemmas:

\begin{lemma}\label{lemma:same_prefrence} Within the Plackett-Luce framework, including the special case of Bradley-Terry, any pair of reward functions belonging to the same equivalence class yield identical preference distributions.
\end{lemma}

\begin{lemma}\label{lemma:same_policy}
    Any two reward functions from the same equivalence class will lead to the same optimal policy in the corresponding constrained reinforcement learning (RL) setup.
\end{lemma}

The arguments for these lemmas are elementary, and details can be found in Appendix \ref{app:lemma1}. The first lemma reflects a well-documented under-specification inherent in the Plackett-Luce class of models \cite{plackett1975analysis}. This ambiguity motivates the necessity for extra identifiability conditions when aiming to establish guarantees about the maximum likelihood estimates from Eq. \ref{eq:reward_model} \cite{bong2022generalized}. The second lemma establishes that the optimal policy is invariant within each reward equivalence class, implying that for practical purposes, it suffices to identify a representative reward function from the target equivalence class. We establish the following theorem in Appendix~\ref{app:thm1}:

\begin{theorem}\label{thm:main}
    Provided mild regularity conditions, all reward equivalence classes that conform to the Plackett-Luce (including Bradley-Terry) models can be described via the reparameterization ${r(x, y) = \beta \log \frac{\pi(y\mid x)}{\pi_\text{ref}(y\mid x)}}$ for some choice of model $\pi(y\mid x)$ and fixed reference model $\pi_\text{ref}(y \mid x)$.
\end{theorem}

\begin{sproof}
    Let $r(x, y)$ be any reward function which induces the optimal policy $\pi_r(y \mid x)$ as specified in Eq. \ref{eq:op_policy}. We intend to show that some reward function within the equivalence class of $r$ admits the specified reparameterized form. Define the projection operator $f$ as follows:  
\begin{equation}
    f(r; \pi_\text{ref}, \beta)(x, y) = r(x, y) - \beta\log\sum_{y}\pi_\text{ref}(y\mid x)\exp\left(\frac{1}{\beta}r(x, y)\right)
\end{equation}
This operation renormalizes $r$ by subtracting the logarithm of the partition function corresponding to $\pi_r$; crucially, this normalization is dependent only on $x$, making $f(r; \pi_\text{ref}, \beta)(x, y)$ another reward function in the equivalence class of $r(x, y)$. Substituting $r$ on the right side of Eq.~\ref{eq:main_eq} (which is valid for any $r$), one obtains $f(r; \pi_\text{ref}, \beta)(x, y) = \beta \log \frac{\pi_r(y\mid x)}{\pi_\text{ref}(y\mid x)}$. Thus, $f$ maps $r$ to a canonical representative within its equivalence class, matching the desired structural form, confirming that our proposed reparameterization incurs no loss of generality in representing reward functions.
\end{sproof}

Theorem~\ref{thm:main} can also be interpreted as pinpointing the precise reward function, within each equivalence class, that is chosen by the DPO reparameterization. Specifically, it selects the reward function that fulfills the following criterion:

\begin{equation}\label{eq:lag_p}
     \sum_{y}\underbrace{\pi_\text{ref}(y\mid x)\exp\left(\frac{1}{\beta}r(x, y)\right)}_{=\pi(y\mid x)\text{, using Thm.~\ref{thm:main} reparam.}} = 1,
\end{equation}

meaning that $\pi(y\mid x)$ constitutes a legitimate probability distribution (all probabilities are non-negative and sum to 1). Additionally, as established by Eq.~\ref{eq:op_policy}, we observe that Eq.~\ref{eq:lag_p} defines the partition function associated with the optimal policy dictated by the reward function $r(x, y)$. The central observation behind the DPO approach is that one can apply specific constraints to the generally under-determined Plackett-Luce (and especially Bradley-Terry) preference model family. This procedure maintains the set of achievable reward models, yet it explicitly ensures the analytical tractability of the optimal policy in Eq.~\ref{eq:op_policy} across all prompt inputs $x$.

\subsection{Instability of Actor-Critic Algorithms} Our framework also enables the identification of potential instability issues present in standard actor-critic algorithms utilized within RLHF, such as PPO. We consider the RLHF process with particular emphasis on the RL fine-tuning component described in Section \ref{section:prelims}. By leveraging connections to the control-as-inference formalism \cite{levine2018reinforcement} for the constrained RL task from \ref{eq:RL}, we proceed with a parameterized policy model $\pi_{\theta}(y\mid x)$ and aim to minimize $\mathbb{D}_{\text{KL}}[\pi_{\theta}(y|x) \mid \mid \pi^*(y\mid x)]$, where $\pi^*$ represents the optimal policy derived in Eq.~\ref{eq:optimum_model} under the reward specification $r_{\phi}(y, x)$. After some mathematical manipulations, this yields the following maximization objective:

\begin{equation}\label{eq:AC}
    \max_{\pi_{\theta}}\mathbb{E}_{\pi_{\theta}(y\mid x)}\bigg[\underbrace{r_{\phi}(x, y) -\beta\log\sum_{y}\pi_\text{ref}(y\mid x)\exp\left(\frac{1}{\beta}r_{\phi}(x, y)\right)}_{f(r_{\phi}, \pi_\text{ref}, \beta)} - \underbrace{\beta\log\frac{\pi_{\theta}(y\mid x)}{\pi_\text{ref}(y\mid x)}}_{\text{KL}}\bigg]
\end{equation}

This objective is identical to that which has been optimized in previous research 
\citep{ziegler2020finetuning, stiennon2022learning, bai2022training, ouyang2022training}, utilizing the DPO-equivalent reward associated with the reward family $r_{\phi}$. Within this context, the normalization factor present in $f(r_{\phi}, \pi_\text{ref}, \beta)$ can be interpreted as the soft value function for the reference policy $\pi_\text{ref}$. Although this normalization component does not alter the optimality of the solution, omitting it can lead to high variance in the policy gradient estimates, which in turn may cause unstable training dynamics. One approach to account for the normalization factor involves learning a value function; however, this is often challenging to optimize in practice. Historically, other works have normalized the rewards by leveraging a human completion baseline, which is effectively a single-sample Monte Carlo approximation of the normalizer. By contrast, the DPO reparameterization introduces a reward function formulation that eliminates the need for such baselines altogether.

\section{Experiments}
Here, we conduct empirical studies to assess DPO's capability to directly train policies from preference data. We start with controlled experiments in the domain of text generation to investigate how effectively DPO balances reward maximization against KL-divergence minimization with the reference policy, compared to standard preference optimization methods such as PPO. Subsequently, we expand our evaluation to encompass larger models and more challenging RLHF benchmarks, including both summarization and dialogue tasks. Our findings indicate that, even with minimal hyperparameter tuning, DPO generally matches or surpasses robust baselines like RLHF implemented with PPO, as well as approaches based on selecting the best of $N$ trajectories according to a learned reward function. Prior to presenting these empirical results, we outline the design of our experimental protocols; further specifics can be found in Appendix~\ref{app:exp_details}.

\textbf{Tasks.} Our study investigates three distinct open-ended text generation scenarios. Across all settings, learning algorithms derive a policy based on a collection of preferences $\mathcal{D}=\bigl\{x^{(i)}, y_w^{(i)}, y_l^{(i)}\bigr\}_{i=1}^N$. For \textbf{controlled sentiment generation}, $x$ represents the beginning of a movie review sampled from the IMDb dataset \cite{maas-EtAl:2011:ACL-HLT2011}, and the policy’s objective is to generate a continuation $y$ that expresses positive sentiment. To facilitate a controlled evaluation, we synthetically construct preference pairs over generated outputs by employing a pre-trained sentiment classifier, ensuring that $p(\text{positive}\mid x,y_w)>p(\text{positive}\mid x,y_l)$. In the case of SFT, GPT-2-large is further trained to convergence using movie reviews from the IMDb dataset’s training partition (additional details in App~\ref{app:sentiment_details}). For the \textbf{summarization} task, $x$ consists of a Reddit forum post, with the policy required to summarize the principal content into $y$. Adopting prior methodologies, we utilize the Reddit TL;DR summarization corpus \citep{volske-etal-2017-tl} in tandem with the set of human preferences collected by \citeauthor{stiennon2022learning}. Here, the SFT model is fine-tuned on human-generated summaries of forum posts\footnote{\url{https://huggingface.co/CarperAI/openai_summarize_tldr_sft}}, and the RLHF experiments are conducted using the TRLX \citep{leandro_von_werra_2023_7790115} toolkit. The human preference annotations come from \citeauthor{stiennon2022learning}, who labeled outputs from a related but separately fine-tuned SFT model. Lastly, for the \textbf{single-turn dialogue} setting, $x$ is an arbitrary user prompt—ranging from astrophysical inquiries to requests for personal advice. The policy must then generate a relevant, informative, and friendly answer $y$. This experiment relies on the Anthropic Helpful and Harmless dialogue dataset \citep{bai2022training}, which consists of 170,000 exchanges between a human user and an AI assistant. Each dialogue concludes with two candidate replies from a large (undisclosed) language model, together with a human annotation of the preferred response. Since an SFT model pre-trained for this domain is unavailable, we construct one by additional fine-tuning of a general language model solely on those completions identified as preferred.

\textbf{Evaluation.} We adopt two complementary strategies to assess model performance in our experiments. In the controlled sentiment generation scenario, where the true reward function (i.e., a sentiment classifier) is available, we systematically evaluate each method based on the trade-off frontier between reward attainment and KL-divergence from the baseline policy. This comparison is possible due to direct access to the ground-truth metric. In contrast, in realistic settings where the underlying reward function is unknown, we measure model quality using the \textit{win rate} against a reference policy, employing GPT-4 as an automated surrogate for human evaluation. For summarization tasks, the reference is the test set's gold summaries, while for dialogue, it is the human-preferred reply from the test data. Although prior work indicates language models can serve as more reliable automatic evaluators than standard metrics \citep{Chen2023ExploringTU}, we bolster our choice of GPT-4 with an independent human study, presented in Sec.~\ref{sec:human-judgments}. Our findings show GPT-4's evaluation outcomes exhibit high alignment with human raters, with human–GPT-4 agreement rates equal to or exceeding inter-annotator consensus among human judges.

\textbf{Methods.} Beyond DPO, we benchmark multiple established approaches for aligning language models with human preferences. As baseline strategies, we include zero-shot prompting using \textbf{GPT-J} \citep{gpt-j} on summarization tasks, and two-shot prompting of \textbf{Pythia-2.8B} \citep{biderman2023pythia} for single-turn dialogue. Further, we assess the \textbf{SFT} approach as well as \textbf{Preferred-FT}, which is produced by fine-tuning a model via supervised learning on the preferred output $y_w$, using outputs from SFT (in controlled sentiment and summarization) or a standard LM (in dialogue). We additionally evaluate the pseudo-supervised \textbf{Unlikelihood} method \citep{welleck2019neural}, which explicitly increases the probability of $y_w$ while suppressing the probability of $y_l$, optionally scaling the penalty via a coefficient $\alpha \in [0, 1]$. The \textbf{PPO} method \citep{schulman2017proximal} is also considered, where the reward is learned from human preference data, alongside the \textbf{PPO-GT} oracle variant, which optimizes using access to the ground truth reward in the sentiment-controlled scenario. For PPO-GT in sentiment experiments, we utilize two versions: an off-the-shelf implementation \cite{leandro_von_werra_2023_7790115}, and a customized version with reward normalization and hyperparameter tuning for enhanced results (these adjustments are similarly applied when using PPO with learned rewards). As a strong baseline, we implement the \textbf{Best of $N$} strategy: here, $N$ outputs are sampled from the SFT (or Preferred-FT for dialogue), and the top-scoring sample, as judged by a reward model trained on preference data, is selected. Although highly effective, this approach is computationally costly since it requires generating $N$ completions for every prompt at evaluation time, making it impractical for larger $N$.

\subsection{How well can DPO optimize the RLHF objective?}

In standard RLHF methods, the optimization objective seeks to maximize reward while enforcing a KL-divergence constraint to prevent the policy from straying too far from the reference policy. Consequently, any evaluation of algorithmic performance must jointly consider both the achieved reward and the associated KL divergence; increases in reward that come at the cost of substantially greater KL may not be desirable. In Figure~\ref{fig:frontier-tldr-main}, we plot the trade-off curves between reward and KL for various algorithms under the sentiment task. Each algorithm is trained multiple times, varying the key hyperparameters governing policy conservativeness—specifically, using target KL values $\in\{3,6,9,12\}$ for PPO, $\beta \in \{0.05,0.1,1,5\}$ and $\alpha\in\{0.05,0.1,0.5,1\}$ for unlikelihood, and employing different random seeds for preferred-FT—resulting in a total of 22 experiments. At every 100-step interval during training up to convergence, we assess the policies on a suite of test prompts, calculating both the mean reward with respect to the actual reward function and the mean sequence-level KL,\footnote{Specifically, this is defined as the aggregate of KL-divergences at each timestep.} relative to the reference policy, denoted by $\text{KL}\left(\pi\mid \mid \pi_\text{ref}\right)$. Our results indicate that DPO yields a significantly superior frontier, achieving high reward without incurring large KL divergence. This finding stands out for several reasons: (1) despite DPO and PPO targeting an identical optimization objective, DPO demonstrates markedly higher efficiency, consistently outperforming PPO in terms of the reward/KL trade-off; and (2) DPO’s performance frontier remains superior to PPO’s even in the scenario where PPO has access to exact reward information (PPO-GT).

\subsection{Can DPO scale to real preference datasets?}
\label{sec:dpo-real-datasets}
We proceed by assessing the fine-tuning capabilities of DPO in the contexts of summarization and single-turn dialogue. In summarization tasks, widely-used automated evaluation metrics such as ROUGE often exhibit weak alignment with human judgments~\citep{stiennon2022learning}. Previous research indicates that applying PPO to LMs, guided by human preference data, yields more effective summaries. Our evaluation involves generating completions using several approaches on the test partition of the TL;DR summarization dataset, then calculating the mean win rate compared to reference completions within this set. For all methods, completions are sampled across a temperature range from 0.0 to 1.0, with the win rates depicted in Figure~\ref{fig:frontier-tldr-main} (right panel). DPO, PPO, and Preferred-FT each use the same underlying GPT-J SFT model\footnote{\url{https://huggingface.co/CarperAI/openai_summarize_tldr_sft}} for fine-tuning. Our results show that at a temperature of 0.0, DPO attains a win rate of about 61\%, outperforming PPO, which reaches approximately 57\% at its own optimal temperature. DPO also surpasses the best-of-$N$ baseline in terms of maximum win rate. It is important to mention that DPO’s $\beta$ hyperparameter was not extensively optimized, implying that these outcomes may not reflect DPO’s full capabilities. Additionally, DPO demonstrates significantly greater robustness across different sampling temperatures than PPO, whose performance may drop to that of the unmodified GPT-J model under higher temperatures. Preferred-FT yields minimal gains over the SFT baseline. Further comparison between DPO and PPO is presented in Section~\ref{sec:human-judgments} using human evaluations, where samples from DPO at temperature 0.25 are favored 58\% of the time relative to PPO samples drawn at temperature 0.

For single-turn dialogue evaluation, we assess various approaches using the portion of the Anthropic HH dataset's test split \citep{bai2022training} that contains a single round of interaction between a human and an assistant. To determine win rates for each technique, GPT-4 evaluates completions by comparing them to the preferred responses in the test data. Because no established SFT model exists for this scenario, we initialize with a pre-trained Pythia-2.8B, apply Preferred-FT to train a reference model on the preferred completions (ensuring output distributions match the model’s), and subsequently fine-tune using DPO. Our analysis also benchmarks against both the best outcome from 128 Preferred-FT generations (since we observe that the Best of $N$ metric reaches a plateau at 128 samples; see Appendix Figure~\ref{fig:best-of-n}) and a 2-shot prompted configuration of the base Pythia-2.8B model, finding DPO matches or surpasses the performance of other methods at their optimal temperature settings. In addition, we examine an RLHF model (trained with PPO on the Anthropic HH dataset\footnote{\url{https://huggingface.co/reciprocate/ppo_hh_pythia-6B}} and sourced from a reputable repository\footnote{\url{https://github.com/CarperAI/trlx/tree/main/examples/hh}}), but are unable to identify prompt or sampling settings where its output exceeds the base Pythia-2.8B. Drawing on both TL;DR findings and the shared reward function between methods, we treat the Best of 128 as a general indicator of PPO-level performance. In summary, DPO is unique among efficient approaches for surpassing preferred completions on the Anthropic HH set, offering results on par with or superior to the computationally intensive Best of 128 baseline. Figure~\ref{fig:dialogue-main} further illustrates DPO’s rapid convergence to peak performance.

\subsection{Generalization to a new input distribution}

To further assess how PPO and DPO handle changes in data distribution, we test the trained policies from our Reddit TL;DR summarization experiment on a distinct corpus—namely, news articles from the CNN/DailyMail test set \citep{nallapati-etal-2016-abstractive}. For consistency, we use the optimal sampling temperatures identified on the TL;DR data (0 and 0.25). Table~\ref{tab:ood} summarizes the outcomes. To evaluate summary quality, we calculate GPT-4 win rates relative to reference summaries in these datasets, employing the same GPT-4 (C) prompt previously used for Reddit TL;DR, substituting the phrase ``forum post'' with ``news article''. On this unfamiliar input distribution, DPO substantially outperforms PPO. These findings offer preliminary support for the hypothesis that DPO-trained policies can generalize to new domains at least as effectively as PPO, despite not benefiting from the additional unlabeled Reddit TL;DR prompts leveraged by PPO.

\subsection{Validating GPT-4 judgments with human judgments}
\label{sec:human-judgments}
We implement a human evaluation study to assess the alignment between GPT-4's assessments and human preferences, focusing on outcomes from the TL;DR summarization task and utilizing two distinct prompting strategies. The \textbf{GPT-4 (S)} (simple) prompt instructs the model to judge which summary better encapsulates the critical information from the post. By contrast, the \textbf{GPT-4 (C)} (concise) prompt explicitly directs the model to also consider conciseness, an important distinction since under the simple prompt, GPT-4 exhibits a tendency to favor longer, more repetitive outputs compared to human evaluators. Full prompt wordings are included in Appendix~\ref{app:prompts}. To capture a wide range of sample qualities, we select three models for comparison: the highest-performing (DPO at temperature 0.25), the lowest (PPO at temperature 1.0), and a model of intermediate performance (SFT at temperature 0.25). Each of these is evaluated against the PPO model using greedy sampling (its optimal temperature setting). Our results reveal that, under both prompt variations, GPT-4's preferences align with human judgments at a rate comparable to human inter-rater agreement, indicating that GPT-4 serves as a credible surrogate for human assessment (although, due to limited annotator availability, we collect repeated human ratings only for DPO and PPO-1 comparisons). Notably, the \textbf{GPT-4 (C)} prompt yields win rates more reflective of human preferences; therefore, we adopt this prompt for reporting principal results in Section~\ref{sec:dpo-real-datasets}. Additional details concerning the setup of the human study—including the rater interface and the roster of participants—are provided in Appendix~\ref{app:human-study}.

\section{Discussion}
Preference-based learning provides a robust and flexible method for developing language models that are both effective and aligned with human values. In this work, we presented DPO, a streamlined approach for training language models from preferences that circumvents the need for reinforcement learning. Instead of reframing the preference learning task to fit the mold of conventional RL so as to leverage established RL techniques, DPO establishes a direct connection between language model policies and reward functions. This approach allows language models to be trained to meet human preferences through a straightforward cross-entropy objective, eliminating the requirement for reinforcement learning without sacrificing generality. DPO matches or surpasses the performance of prevalent RLHF methods, such as those employing PPO, while requiring minimal hyperparameter adjustment; thus, it lowers the complexity associated with preference-based training for language models.

\textbf{Limitations \& Future Work.} The findings in this study suggest multiple avenues for further research. For instance, it remains to be seen how well the DPO-trained policy extrapolates to data outside of its training distribution, particularly in comparison to models optimized via explicit reward functions. Preliminary analyses indicate that the generalization capabilities of DPO are comparable to those observed with PPO-based approaches, yet more thorough investigation is warranted. Additionally, there is an open question as to whether self-labeling strategies based on DPO can efficiently leverage prompts without explicit labels. Another pertinent issue is understanding how reward over-optimization occurs under direct preference optimization and whether the marginal drop in performance depicted in Figure~\ref{fig:dialogue-main}-right exemplifies this effect. While our experiments were conducted on models with up to 6B parameters, scaling DPO to contemporary large-scale architectures presents an intriguing challenge for upcoming studies. Concerning evaluation, we noticed that win rates generated by GPT-4 are influenced by the prompt used; this highlights an area for future exploration regarding optimal prompt design to obtain reliable judgments from automated evaluators. Lastly, the potential applications of DPO extend well beyond language modeling from human preferences, and encompass training generative models across different data modalities.